%%%%%%%%%%%%%%%%%%%%%%%%%%%%%%%%%%%%%%%%%%%%%%%%%%%%%%%%%%%%%%%%%%
\section{Introduction}
\label{sec:intro}
%%%%%%%%%%%%%%%%%%%%%%%%%%%%%%%%%%%%%%%%%%%%%%%%%%%%%%%%%%%%%%%%%%

\noindent
Mooring lines, risers, aquaculture netting, textile yarns, and aquatic vegetation all have a cross-section orders of magnitude smaller than the flow domain they occupy, and they all deform under the flow that loads them. Resolving such a body with a boundary-fitted mesh imposes severe local refinement, and the mesh must then be deformed or regenerated at every time step as the structure moves. The cost becomes prohibitive for a canopy of vegetation stems or a net of many strands.

\noindent
Distributed forcing methods avoid that cost by representing the structure through a lower-dimensional description and introducing its action on the fluid as regularised momentum sources on a fixed Eulerian mesh. The immersed boundary method \citep{peskin_flow_1972,peskin_immersed_2002,mittal_immersed_2005,uhlmann_immersed_2005} established the idea for surfaces, and the actuator line method of \citet{sorensen_numerical_2002} reduced it further for slender members by concentrating the hydrodynamic interaction at actuator points along a curve. The method was developed for wind-turbine wakes. \citet{troldborg_numerical_2010} showed that it reproduces the essential momentum exchange on meshes far coarser than blade-resolved calculations. It has since been refined in three directions, namely the velocity sampled to drive the blade force \citep{churchfield_advanced_2017}, the shape and width of the smoothing kernel \citep{jha_actuator_2018}, and the generalisation of the actuator curve away from a straight radial line \citep{navarro_diaz_actuator_2023}.

\noindent
Nothing in the concept restricts it to rotating blades. \citet{bral_development_2025} adapted the actuator line to flexible yarns in air-jet weaving, and introduced the upstream sampling strategy discussed below to keep the drag estimate away from the momentum deficit the line itself creates. \citet{bral_modeling_2025} then extended that formulation into a two-way coupled framework in which the sampled load drives a beam model and the beam displacement feeds back into the source term. \citet{schmitt_coupled_2022} took a different route to the same class of problem, coupling an actuator line in OpenFOAM to an external finite element structural solver, and applied it to flexible turbine blades and fish-farm nets. These three studies establish that a drag-loaded actuator line can drive a deforming slender structure, and they are the closest antecedents of the present work.

\noindent
Two operators govern the accuracy of such a coupling. The first is velocity sampling, for which the usual choices are point, integral, and line-average strategies \citep{sorensen_numerical_2002,merabet_parametric_2019,zormpa_effect_2025}. The distinction matters most for drag-dominated members, whose load is aligned with the relative velocity, so that the regularised source acts as a local sink slowing the fluid in the same region where the velocity would be sampled. Because the drag scales with the square of that velocity, an underestimate feeds back on itself. \citet{bral_development_2025} broke this feedback by sampling at a fixed distance upstream of the actuator point, outside the momentum-deficit region, and the same concern underlies the sampling and smearing corrections of \citet{churchfield_modeling_2015}, \citet{martinez_tossas_optimal_2017} and \citet{muscari_effective_2024}. The second operator is force regularisation, since a point force cannot be applied on a finite volume mesh without severe grid dependence and is instead spread with a Gaussian kernel of width $\varepsilon$. A kernel that is too narrow makes the source grid-localised, and one that is too wide weakens the local effect of the line \citep{martinez_tossas_optimal_2017,navarro_diaz_actuator_2023}. Beyond accuracy, the two operators determine which discrete properties the coupling preserves, namely the conservation of exchanged momentum, the resolution of the kernel relative to the local grid, and invariance under parallel domain decomposition \citep{bral_development_2025}.

\noindent
On the structural side, the geometrically exact beam theory of \citet{reissner_finite_1981}, extended to three dimensions by Simo, treats finite rotations without linearisation and retains shear and torsion alongside bending and axial stretch. Most implementations of it are finite element, and numerical studies of flexible vegetation illustrate what follows from that choice. \citet{chen_eulerianlagrangian_2019} coupled a RANS--VOF flow solver in OpenFOAM to a finite element elastic rod through a diffused immersed-boundary force, and showed that large vegetation deflections can be reproduced without a body-fitted mesh, validating against both single-stem and patch experiments. \citet{elrahi_numerical_2023} coupled a smoothed-particle flow solver to an external finite element library for the same class of problem, and \citet{yin_experimental_2023} combined experiments with numerical modelling of emergent flexible vegetation under waves. \citet{kim_steady_2024} used an immersed-boundary soft-body formulation to study steady flow over a finite patch, and \citet{pruter_comprehensive_2025} applied a direct-forcing immersed-boundary model to current-induced vegetation interaction. \citet{dickinson_modelling_2025} represented the stem instead as an elastically articulated body.

\noindent
In every one of these studies the structure is described by a different numerical method from the flow. The two sides therefore rest on different assumptions about how the governing equations are integrated, hold their data in different structures, and are decomposed differently across processors, and each must be verified in its own setting before the coupled result can be relied upon. A cell-centred finite volume discretisation of the same beam theory \citep{bali_cell-centered_2022,bali_finite_2024,bali_beamfoam_2025} removes that separation, because the two subsystems then share one set of conservation statements, one mesh and field infrastructure, one domain decomposition, and one time loop. Two-way actuator-line coupling has been demonstrated for flexible marine structures, but on finite element or external structural solvers \citep{schmitt_coupled_2022,bral_modeling_2025}, so the sampling, regularisation, and parallel exchange have not been established for a geometrically exact finite volume beam.

\noindent
The application that motivates a line-based representation is the canopy rather than the single stem, since a meadow contains many stems and none of them can carry a body-fitted mesh at that scale. Two strategies are visible in the literature. \citet{kim_steady_2024} resolved a finite patch member by member, which captures the flow between neighbours but inherits the per-stem cost, whereas \citet{jin_numerical_2025} moved in the opposite direction and represented vegetation flows semi-resolved, through distributed virtual body forces that stand in for the stems. An actuator-line beam sits between the two, keeping an individual deformable structure for every stem while paying only the cost of a line. The per-stem accuracy of such a model is what a canopy calculation inherits, which is why it is quantified here on isolated stems before any canopy is attempted.

\noindent
The hydrodynamic closure adopted here is a quasi-steady Morison-type drag law with no added-mass or history term, and a closure of this kind has an applicability boundary that can be stated in advance. A complementary experimental and theoretical literature supplies that boundary. \citet{luhar_waveinduced_2016} measured the wave-induced motion of individual flexible blades and organised the response through the Cauchy number and the ratio of blade length to wave orbital excursion, describing the drag reduction produced by reconfiguration as an effective length. \citet{leclercq_reconfiguration_2018} then showed that the character of the response is governed by the combination $CaL/KC$, which is proportional to the squared ratio of the natural period of the fluid-loaded blade to the wave period. Below a threshold of approximately $0.2$ the blade responds quasi-statically in a single spatial mode, and above it the added mass of the surrounding water amplifies the motion and excites higher modes. \citet{jacobsen_waveinduced_2019} confirmed this classification experimentally for single vegetation mimics, measuring the motion along the entire stem together with the base force, and derived drag coefficients by fitting a relative-velocity Morison formulation to that force. They reported that the fitted coefficients decrease with the Keulegan--Carpenter number, carry a documented spread between repeat tests, and reproduce the measured forces generally within $10$ to $20\%$, while the reconstructed inertia force contributes less than $10\%$ of the maximum force for quasi-static blades.

\noindent
Published drag coefficients therefore carry an irreducible band, so a study that fixes the coefficient in advance should propagate that band through the simulation rather than tune a single value against the response being validated. A quasi-steady drag-only closure also has an a priori boundary at the quasi-static regime limit, beyond which added-mass inertia, and in principle the history forces of the Morison family \citep{sarpkaya_mechanics_1981}, become dynamically significant and such a closure must be expected to under-predict.

\noindent
This paper formulates and implements a two-way actuator-line coupling between a cell-centred finite volume Simo--Reissner beam and an incompressible finite volume flow solver, and verifies and validates it against three reference problems of increasing difficulty. The coupling requires no external structural code, because the beam shares the mesh, fields, time loop, and domain decomposition of the flow solver. Its sampling and projection operators are parallel consistent on a domain-decomposed fluid mesh, and the normalisation applied to the projection makes the transferred resultant force exact to round-off independently of the mesh, the kernel width, and the kernel truncation. The formulation does not depend on the number of fluid phases, and it is used unchanged in single-phase channel flow and in a two-phase wave flume. Every drag coefficient used in the verification and validation campaign is fixed in advance. That campaign also identifies the two operator settings with a physical rather than a purely numerical influence on the response, and it quantifies the accuracy of the closure together with the range of conditions over which that accuracy holds. Verification and validation are used throughout in the standard sense \citep{roache_quantification_1997,oberkampf_roy_2010}, the first asking whether the coupled equations are solved correctly and established against independent numerical solutions, the second asking whether the model reproduces physical reality and established against experiment.

\noindent
Section~\ref{sec:mathmodel} presents the governing equations and the hydrodynamic closure, and Section~\ref{sec:numericalmodel} the finite volume beam discretisation and the coupling operators. Sections~\ref{sec:cantilever} to~\ref{sec:leaflet} verify the structural solver and then the coupling, against a flexible beam in a steady confined channel and a leaflet in a sinusoidal channel flow. Section~\ref{sec:flexveg} validates the coupling against the wave-flume experiment of \citet{jacobsen_waveinduced_2019} on a flexible vegetation stem, and Section~\ref{sec:conclusion} draws the conclusions.
